% Scheletro di lavoro, non testo definitivo.
% Fonti e passi sono proposte di lettura: verificarli sul testo greco prima della stesura.
% Seguira il metodo: documento originale, traduzione, esame dell'argomentazione, ricostruzione critica.
\chapter{Il dialogo filosofico: scrittura, cronologia ed esegesi}
\label{cap:platone:dialogo}

% Problema: Il protagonista di un dialogo esprime necessariamente il pensiero di Platone?
% Fonti da verificare: Apologia; Eutifrone; Repubblica I.
\section{Il problema dell'interpretazione}

\subsection{L'autore, i personaggi e l'assenza della voce narrante}

\subsection{Aporia, confutazione e ricerca della definizione}

% Problema: Che cosa aggiungono scena, interlocutori e miti alla struttura argomentativa?
% Fonti da verificare: Protagora; Gorgia; Fedro.
\section{La forma del dialogo}

\subsection{La situazione drammatica e la costruzione dell'argomento}

\subsection{Mito, immagini e discorso razionale}

% Problema: Come interpretare la critica della scrittura e le testimonianze sull'oralità?
% Fonti da verificare: Fedro 274b–278e; Lettera VII 341c–345c (autenticità controversa).
\section{Scrittura e insegnamento}

\subsection{La critica della scrittura nel Fedro}

\subsection{La Lettera VII e la questione delle dottrine non scritte}

% Problema: In quale misura la datazione delle opere può orientare l'esegesi?
% Fonti da verificare: Testi dialogici; studi stilometrici e testimonianze antiche da individuare.
\section{La cronologia dei dialoghi}

\subsection{Criteri, limiti e ipotesi di datazione}

\subsection{Lettura genetica e lettura problematica}
